\section{Introduction}

In modern information systems, knowledge is gathered from vast quantities of various sources. These systems form the basis of many tools that we use today, ranging from search engines and knowledge bases to recommendation systems and scientific databases. One of the largest repositories of information used by these systems is the web. Today, the web is experiencing a fundamental shift from human-to-human to machine-to-machine information exchange.

With advancements in artificial intelligence systems, particularly large language models (LLMs) and autonomous agents, this transition has accelerated rapidly. Information is now being generated, distributed, and consumed with minimal or no human intervention. Agents routinely retrieve information from several independent sources and use the result to make decisions and execute downstream tasks autonomously. Hence, the credibility of this information is critical for the actions that follow. Existing approaches to validate such information typically rely on external evaluation such as human review or verification against trusted sources. More recently, information is often routed through a secondary large language model for evaluation. These approaches can be effective, but constantly relying on external evaluation for every decision can be slow and computationally expensive while becoming increasingly difficult to integrate as autonomous systems continue to operate at greater pace and scale. This is because credibility cannot be computed within the information network itself, but instead must be established through external evaluation.

To bridge this structural gap, we introduce Verity, a graph-based framework for inferring source credibility in large collections of conflicting claims. It estimates source credibility from the network of assertions connecting sources and claims. These credibility estimates allow autonomous systems to distinguish between conflicting information while avoiding reliance on external evaluation.
